\chapter*{Anexo C: Modelo de datos de la plataforma}
\addcontentsline{toc}{chapter}{Anexo C: Modelo de datos de la plataforma}
\label{cap:anexo_modelo}

\setcounter{section}{0}
\renewcommand{\thesection}{C.\arabic{section}}

Este anexo describe las estructuras de datos que atraviesan la plataforma,
desde el sobre de procedencia de bronze hasta las tablas publicadas en
PostgreSQL. Las definiciones proceden del código del repositorio, de los
ficheros \path{src/tfm_lakehouse/silver/schemas.py} y
\path{src/tfm_lakehouse/gold/schemas.py}, que las declaran como datos y no
como tipos de Spark para poder comprobarlas sin ejecutar Spark.

\section{Organización del almacenamiento}
\label{sec:anexo_modelo_almacenamiento}

Bronze guarda ficheros de texto en formato JSON Lines bajo el prefijo
\path{bronze/} del cubo \texttt{lakehouse}. Cada fuente tiene su propia
ruta, y cada ejecución escribe en una partición con su fecha y su
identificador, con un fichero de manifiesto que marca la partición como
completa:

\begin{center}
\small\path{bronze/source=<fuente>/ingestion_date=<fecha>/run_id=<ejecución>/part-NNNNN.jsonl}
\end{center}

Los registros que no superan la comprobación de aterrizaje se guardan en
una ruta paralela de rechazados con los mismos niveles. Silver y gold son
tablas Iceberg bajo la raíz \path{warehouse/} del mismo cubo, en los espacios
de nombres \path{lakehouse.silver} y \path{lakehouse.gold}.

\section{Sobre de procedencia de bronze}
\label{sec:anexo_modelo_sobre}

Cada línea aterrizada es un sobre con los campos de la
Tabla~\ref{tab:anexo_sobre}. El contenido original de la fuente viaja sin
modificar en el campo de carga útil, como texto XML para ORCID masivo y como
objeto JSON para la API y para el CVN sintético.

\begin{table}[H]
  \centering
  \footnotesize
  \renewcommand{\arraystretch}{1.1}
  \begin{tabular}{
    >{\raggedright\arraybackslash}p{0.3\textwidth}
    >{\raggedright\arraybackslash}p{0.58\textwidth}}
    \hline
    \textbf{Campo} & \textbf{Contenido} \\
    \hline
    \texttt{record\_id} & Identificador único del registro, base de la deduplicación. \\
    \texttt{source} & Fuente de origen: \texttt{orcid\_bulk}, \texttt{orcid\_api} o \texttt{synthetic\_cvn}. \\
    \texttt{source\_ref} & Referencia dentro de la fuente, como un identificador ORCID o un identificador de documento CVN. \\
    \texttt{source\_snapshot} & Identificador de la instantánea de origen, que incluye el límite de registros aterrizados. \\
    \texttt{source\_files}, \texttt{source\_artifacts} & Ficheros y artefactos de origen, con el vocabulario del bloque de trazabilidad de Open CVN. \\
    \texttt{retrieved\_at} & Momento en que se obtuvo el dato de su fuente. \\
    \texttt{landed\_at}, \texttt{ingestion\_date} & Momento y fecha de la ejecución que lo aterrizó. \\
    \texttt{ingestion\_run\_id} & Identificador de la ejecución, válido como segmento de ruta. \\
    \texttt{landing\_check} & Estado de la comprobación de aterrizaje y sus errores. \\
    \texttt{payload\_format} & Formato de la carga útil, \texttt{xml} o \texttt{json}. \\
    \texttt{payload} & Contenido original de la fuente. \\
    \hline
  \end{tabular}
  \caption{Campos del sobre de procedencia de bronze.}
  \label{tab:anexo_sobre}
\end{table}

\section{Tablas de silver}
\label{sec:anexo_modelo_silver}

Silver tiene seis tablas Iceberg. Todas se reconstruyen por completo en cada
ejecución, y dos ejecuciones sobre la misma entrada producen el mismo
contenido, incluidos los identificadores de entidad. La
Tabla~\ref{tab:anexo_silver} resume sus columnas. Las columnas de
procedencia de \texttt{person\_record} y de \texttt{rejected} copian las del
sobre.

\begin{table}[H]
  \centering
  \footnotesize
  \renewcommand{\arraystretch}{1.1}
  \begin{tabular}{
    >{\raggedright\arraybackslash}p{0.19\textwidth}
    >{\raggedright\arraybackslash}p{0.69\textwidth}}
    \hline
    \textbf{Tabla} & \textbf{Columnas principales} \\
    \hline
    \texttt{person\_record} & Una fila por registro válido de cualquier fuente: \texttt{record\_id}, \texttt{source}, \texttt{source\_ref}, nombres y apellidos con sus formas normalizadas, clave de bloqueo, inicial del nombre, \texttt{orcid\_id}, fecha de última modificación, columnas de procedencia y avisos de validación. \\
    \texttt{affiliation} & Una fila por empleo o formación: tipo, organización y su forma normalizada, cargo, departamento, ciudad, país y año y mes de inicio y de fin. \\
    \texttt{publication} & Una fila por obra: título y su forma normalizada, año, mes, DOI, tipo de obra y autores, que solo existen en los CVN. \\
    \texttt{entity\_link} & Una fila por registro: entidad asignada, regla de fusión, evidencia en JSON y los indicadores de ambigüedad y de conflicto de nombre. \\
    \texttt{entity} & Una fila por persona resuelta: identificador, identificador ORCID, nombre, número de registros, fuentes, reglas usadas y si tiene CVN y ORCID. \\
    \texttt{rejected} & Una fila por registro rechazado: reglas incumplidas, errores y columnas de procedencia. \\
    \hline
  \end{tabular}
  \caption{Tablas de silver.}
  \label{tab:anexo_silver}
\end{table}

Los registros rechazados indican siempre la regla concreta que incumplieron.
Las reglas son once, seis propias del CVN y cinco de ORCID:

\begin{itemize}
  \item \textbf{CVN:} análisis del JSON, esquema del documento, esquema de cada
  entidad, dígito de control del identificador ORCID declarado, presencia de
  los campos obligatorios y orden de las fechas.
  \item \textbf{ORCID:} análisis del XML, forma del JSON de la API, dígito de
  control del identificador, presencia de los campos obligatorios y
  existencia de alguna actividad.
\end{itemize}

Una línea que no es un sobre válido se rechaza antes con una regla propia.

\section{Tablas de gold y su espejo en PostgreSQL}
\label{sec:anexo_modelo_gold}

Gold tiene cinco tablas, escritas en Iceberg y copiadas una a una al esquema
\texttt{gold} de PostgreSQL. La Tabla~\ref{tab:anexo_gold} las describe junto
con sus claves y sus índices secundarios. El mapeo de tipos es directo, y
solo los instantes de tiempo pasan a un tipo con zona horaria.

\begin{table}[H]
  \centering
  \footnotesize
  \renewcommand{\arraystretch}{1.1}
  \begin{tabular}{
    >{\raggedright\arraybackslash}p{0.25\textwidth}
    >{\raggedright\arraybackslash}p{0.34\textwidth}
    >{\raggedright\arraybackslash}p{0.2\textwidth}}
    \hline
    \textbf{Tabla} & \textbf{Contenido} & \textbf{Clave e índices} \\
    \hline
    \path{dim_researcher} & Una fila por entidad, con nombre, fuentes, número de obras distintas, primer y último año de publicación y rango de años de carrera. & Clave por entidad \\
    \path{publications_per_researcher_year} & Número de obras distintas por entidad y año. & Clave por entidad y año, índice por año \\
    \path{affiliation_timeline} & Periodos de afiliación por entidad, con organización, cargo y años. & Índices por entidad y por organización \\
    \path{collaboration_pairs} & Pares de entidades con obras compartidas, años extremos y marca de miembro con CVN. & Clave por par, índice por segunda entidad \\
    \path{gold_run} & Una fila por ejecución, con sus recuentos y las instantáneas de silver leídas. & Clave por ejecución \\
    \hline
  \end{tabular}
  \caption{Tablas de gold publicadas en PostgreSQL.}
  \label{tab:anexo_gold}
\end{table}

La tabla de ejecución se escribe la última, y la publicación se niega a
funcionar si no existe. Su presencia marca así una escritura completa y
permite atribuir cada cifra publicada a las instantáneas exactas de silver
de las que procede. En la ejecución de referencia las cinco tablas tenían
29\,767, 172\,861, 94\,431, 9\,540 y 1 filas, en el orden de la tabla.
